% =============================================================================
% 5.6 Limitations
% =============================================================================

\section{Limitations}
\label{sec:limitations}

\paragraph{Infrastructure representativeness.}
The Kubernetes cluster was a single-node Kind deployment running on the same macOS host as the VM.
This means the Kind cluster and the VM were competing for the same physical CPU and RAM.
A production Kubernetes cluster would have dedicated nodes with separate resource pools, potentially shifting the crossover point.
Similarly, Kind's use of a single node eliminates true multi-node scheduling decisions, pod affinity/anti-affinity logic, and inter-node network overhead.

\paragraph{Memory asymmetry between environments.}
The VM was configured with 4\,GB RAM while the Kind node had 8\,GB.
This 2$\times$ memory advantage directly benefits Kubernetes: the VM's OOM kills at L4+ are driven by aggregate memory demand exceeding the 4\,GB host capacity, while the Kind node's 8\,GB provides sufficient headroom for all concurrency levels.
This asymmetry means the observed crossover point (L4) partially reflects the hardware difference rather than being solely attributable to architectural isolation.
A fairer comparison would match both environments at the same total memory; however, the architectural lesson---that per-pod resource limits prevent cross-job interference regardless of total capacity---remains valid.

\paragraph{Synthetic workload.}
The benchmark workload (SHA-256 hashing + NumPy memory allocation) was designed to generate controlled, reproducible load.
Real Dagster pipelines include external dependencies (database queries, API calls, file system reads) whose latency characteristics are fundamentally different from CPU-bound workloads.
The crossover point may differ for I/O-bound, memory-bound, or network-bound workloads.

\paragraph{Statistical power.}
With three repetitions per concurrency level, the study has limited statistical power to detect small effects.
Results should be interpreted as directional rather than definitive.
The reliability failure rates (66.7\% at L4 through 0.0\% at L6) are large effects with clear statistical significance. The startup overhead range (19--57\,s) grows with concurrency and would require additional investigation under colder image conditions or in production clusters to assess variance.

\paragraph{Reliability threshold exceeded early.}
The formal reliability threshold (success rate $<$95\%) was crossed at L4, earlier than hypothesised.
This is the primary finding, but it also means the crossover analysis is dominated by the reliability dimension rather than the performance dimension.
Future work with a VM having more RAM (8--16\,GB) would permit exploration of the performance-only crossover regime.

\paragraph{Dagster-specific startup overhead.}
Container startup time (19--57\,s in this local Kind environment, growing with concurrency) is specific to the combination of pre-pulled images, in-cluster networking, and the Dagster framework's initialisation sequence. In a production cluster with cold image pulls, startup overhead would be higher (30--90\,s depending on image size and registry speed).
Other workflow frameworks with lighter runtime footprints may have substantially lower startup overhead, which would shift the performance crossover point to lower concurrency levels.

\paragraph{Infrastructure configuration changes during experiments.}
Two fixes were applied between initial deployment and the final experiments: the \texttt{dagster-docker} package was added to the VM daemon environment, and container memory limits were increased.
The final experimental data was collected after both fixes were in place, but the existence of these issues during development means other undetected inconsistencies may remain.

